\documentclass[12pt, a4paper]{ntut-report}
\usepackage[dvips,xetex]{graphicx}
\usepackage{ifpdf,mla}% <-- mla.sty requires ifpdf.sty, but (perversely) doesn't load it
%\usepackage{fontspec}
\usepackage{geometry}
\usepackage{lipsum}
\usepackage{xeCJK}
\usepackage{wallpaper}
\usepackage{pdfpages}
\usepackage{indentfirst}
\usepackage{setspace}
\usepackage{hyperref}
\usepackage{zhnumber}
\usepackage{titlesec}
\usepackage{amsmath}
\usepackage[backend=bibtex,style=ieee,sorting=none]{biblatex}
\usepackage{tabularx}
\usepackage{multirow}
\usepackage{listings}
\usepackage{csquotes}
\usepackage{subcaption}
\usepackage{float}
\usepackage[ruled,vlined,linesnumbered]{algorithm2e}

\MakeAutoQuote{“}{”}
\MakeOuterQuote{"}

\captionsetup{justification=centering,singlelinecheck=false,font={normalsize,stretch=2}}
\captionsetup[sub]{justification=centering,singlelinecheck=false,font={normalsize,stretch=2}}
\SetAlgorithmName{演算法}{演算法}{演算法列表}
\SetKwInput{KwInput}{Input}
\SetKwInput{KwOutput}{Output}
\SetKw{KwRet}{return}

\renewcommand{\baselinestretch}{1.5}
\xeCJKsetup{AutoFakeBold=true, AutoFakeSlant=true}
\setCJKmainfont{DFKai-SB} % 中文字體
\setmainfont{Times New Roman} % 英文字體
\geometry{a4paper,total={297mm,210mm},top=4.5cm,bottom=0.75cm,left=2.5cm,right=2.5cm} % 頁面設定，不需修改
\pagestyle{plain} % 頁碼
\addbibresource{reference.bib}

% 中文碩士論文：作者，論文名稱，文章種類，所屬學校，年份。
\DeclareFieldFormat[thesis]{title}{#1}
% IEEE 文獻（網址結尾除外）應以句點結束，包含以 DOI 結尾者。
% 頁碼不加入千分位分組空格。
\DeclareFieldFormat{pages}{\mkpageprefix[bookpagination]{#1}}
\DeclareFieldFormat{doi}{%
  \mkbibacro{DOI}\addcolon\space
  \ifhyperref
    {\href{https://doi.org/#1}{\nolinkurl{#1}}}
    {\nolinkurl{#1}}}
\DeclareBibliographyDriver{thesis}{%
  \usebibmacro{bibindex}%
  \usebibmacro{begentry}%
  \printnames{author}%
  \setunit{\printtext{，}}%
  \usebibmacro{title}%
  \setunit{\printtext{，}}%
  \printfield{type}%
  \setunit{\printtext{，}}%
  \printlist{institution}%
  \setunit{\printtext{，}}%
  \printfield{year}%
  \begingroup
  \renewcommand*{\finentrypunct}{。}%
  \usebibmacro{finentry}%
  \endgroup}

% 線上文獻：網址前統一標示 [Online; accessed DD-Month-YYYY]。
\newcommand*{\NTUTAccessMonth}[1]{%
  \ifcase#1\relax
  \or January\or February\or March\or April\or May\or June%
  \or July\or August\or September\or October\or November\or December%
  \fi}
\DeclareFieldFormat{url}{Available\addcolon\space\url{#1}\nopunct}
\renewbibmacro*{url+urldate}{%
  \iffieldundef{url}
    {}
    {\iffieldundef{urlyear}
       {\printtext{[Online]}}
       {\printtext{[Online\addsemicolon\space accessed \thefield{urlday}-%
          \NTUTAccessMonth{\thefield{urlmonth}}-%
          \thefield{urlyear}]}}%
     \setunit{\adddot\space}%
     \printfield{url}}}

\input{ntut-labels.tex}

\begin{document}

% 封面、不用浮水印 Cover without a watermark
\include{page/titlepage.tex}
% 留白頁、不用浮水印 Blank page without a watermark
\newpage
\thispagestyle{empty}
\null
% (預留空白頁)


% 新增浮水印 Add Watermarks
\CenterWallPaper{.64}{ntut-logo.png}

% 封面、要浮水印 Cover with a watermark
\input{page/titlepage.tex}

\pagenumbering{roman}

% 摘要頁（中文）Abstract (Chinese)
% If you don’t have Chinese abstract, please delete this one.
% 中文摘要頁
\begin{ZhAbstract}
    \begin{ZhAbstractItems}
        % 關鍵詞，請自己填，請自己填，多個關鍵字以逗號(、)隔開
        \noindent \text 關鍵詞：大型語言模型代理人、API 任務自動化、原子化經驗、動態局部修復

    \end{ZhAbstractItems}

    \begin{ZhAbstractDescription}
        大型語言模型代理人能透過應用程式介面（Application Programming Interface, API）操作外部系統，但在多步驟任務中，若規劃不完整、資料依賴處理錯誤，或個別步驟執行失敗，後續步驟可能因缺少正確的中間結果而無法完成，進而影響整體任務成果。本研究採用的基礎系統為 API 知識圖譜代理人（API Knowledge Graph Agent, APIKGAgent），其主要透過重新規劃處理執行失敗，尚未整合由成功任務軌跡建立的經驗重用機制與局部計畫修復。因此，本研究探討如何運用既有成功經驗，並針對局部錯誤調整計畫，以提升代理人的任務完成能力。
        本研究以 APIKGAgent 為基礎，提出案例式修復 API 知識圖譜代理人（Case-Based Repair API Knowledge Graph Agent, CBR-APIKGAgent），整合原子化經驗（Atomic Experience）重用與動態局部修復（Dynamic Local Repair）。系統從訓練資料的成功任務軌跡中萃取可跨任務重用的策略單元，並依當前任務或失敗情境檢索相關經驗，提供初始規劃、局部修復與重新規劃參考。當步驟執行失敗時，系統根據錯誤訊息、相依資料與執行結果，透過插入、替換或刪除步驟調整局部計畫；若無法形成有效修正，再回退至重新規劃。
        本研究於 AppWorld 的 168 個多應用程式任務上進行評估。完整系統的任務目標完成率由基礎系統的 37.5\% 提升至 51.8\%，情境目標完成率則由 14.3\% 提升至 35.7\%。消融實驗與案例分析顯示，原子化經驗主要協助代理人掌握資料依賴、目標集合與操作條件，動態局部修復則能處理授權缺失、前置資料不足及局部操作不完整等執行錯誤；兩項機制分別作用於規劃與執行階段，結合後取得最佳整體表現。研究亦發現，系統表現仍會受到檢索經驗與任務限制的適配程度、初始規劃品質，以及局部修復與重新規劃判斷等因素影響。儘管存在上述限制，整體結果仍顯示，細粒度經驗重用與局部計畫修復有助於提升大型語言模型代理人在多步驟 API 任務中的自我修正能力。
    \end{ZhAbstractDescription}
    
\end{ZhAbstract}

% 摘要頁（英文）Abstract (English)
% 英文摘要頁
\begin{EnAbstract}
    \begin{EnAbstractItems}
        % 關鍵詞，請自己填，多個關鍵字以逗號 "," 隔開
        \noindent \text Keyword: Large Language Model Agents, Self-Correction Mechanism, Case-Based Learning, Dynamic Plan Repair

    \end{EnAbstractItems}

    \begin{EnAbstractDescription}
        When executing API tool-use tasks, Large Language Model (LLM) agents often fail due to incomplete planning, missing parameters, authentication errors, improper API selection, or misinterpretation of execution results. Existing frameworks predominantly rely on holistic re-planning upon encountering errors, which tends to waste successfully executed steps and may lead to recurring the same mistakes.

        To address this issue, this research proposes a self-correction mechanism for API task agents that integrates case-based learning with dynamic plan repair. The system generates actionable repair feedback through structured error diagnosis and performs local plan adjustments at the point of failure using a "Rerollout" mechanism, which includes inserting, replacing, or deleting repair steps.

        Furthermore, the system extracts successfully repaired cases into atomic experiences, which are then reused through semantic retrieval during the planning, re-planning, and local repair stages. In terms of implementation, this study also incorporates structured plan representation, dependency correction, token synchronization, and observability metric logging to support subsequent experimental analysis.

        Preliminary results indicate that this mechanism provides a more fine-grained error-handling workflow and establishes a quantifiable foundation for evaluating the error recovery capabilities and the effectiveness of experience reuse in agents.    \end{EnAbstractDescription}
    
\end{EnAbstract}


% 鳴謝 Acknowledgements
\begin{Thanks}
    回顧這段碩士求學歷程，從最初的摸索到最終完成論文，過程中經歷了無數次嘗試與修正。能夠走到這一步，並非憑藉一人之力，而是有賴許多人的陪伴與支持。

    首先，誠摯感謝指導教授劉建宏老師。在研究過程中，老師總能在關鍵時刻點出問題核心，引導我重新思考研究方向。無論是研究上的討論，或是面對困難時的提醒，都讓我逐漸建立更成熟的思考方式，也學會以更嚴謹的態度面對學術問題。

    同時，也誠摯感謝梁文耀教授的指導與建議。老師在討論中提出的觀點，經常使我跳脫原有框架，讓研究內容更加周全，也使我對自己的研究有了更深入的理解。

    在實驗室的日子裡，感謝珮倫與郁喬，在我因論文進度感到沮喪、對自己失去信心時，給予我鼓勵與支持；也感謝冠穎與以謙，在我對研究有所疑惑時，耐心提供幫助與建議。此外，也想感謝一路上關心與支持我的朋友們。在我感到疲憊或迷惘時，是你們的鼓勵讓我得以重新調整步伐，繼續向前。

    最後，我想將這份成果獻給我的家人。謝謝你們長久以來的理解與支持，讓我能夠專心投入學業，並在面對挑戰時，始終保有繼續前進的力量。

    這段旅程的結束，同時也是另一段道路的開始。謹以此文，向所有曾經幫助與陪伴我的人，致上最真摯的感謝。

\end{Thanks}

% 目錄、表目錄與圖目錄 Table of Contents, List of tables, List of figures
\begin{TableOfContent}
    \tableofcontents
\end{TableOfContent}

\makeatletter
\begin{TableOfContent}
    \begingroup
        \let\NTUTOriginalNumberline\numberline
        \renewcommand{\numberline}[1]{%
            \NTUTOriginalNumberline{\figurename\enspace #1}%
        }
        \renewcommand*\l@figure{\@dottedtocline{1}{1.5em}{3.8em}}
        \listoffigures
    \endgroup
\end{TableOfContent}

\begin{TableOfContent}
    \begingroup
        \let\NTUTOriginalNumberline\numberline
        \renewcommand{\numberline}[1]{%
            \NTUTOriginalNumberline{\tablename\enspace #1}%
        }
        \renewcommand*\l@table{\@dottedtocline{1}{1.5em}{3.8em}}
        \listoftables
    \endgroup
\end{TableOfContent}
\makeatother

 
\pagenumbering{arabic}

% 章節一 Chapter 1
\chapter{緒論}

\section{研究動機}

本研究聚焦於大型語言模型代理人在多步驟 API 任務中的規劃、執行與錯誤修復能力。當代理人需要跨應用程式完成目標時，單一步驟的失誤可能連鎖影響後續流程，因此有必要引入成功經驗參考與動態修復機制，以提升整體任務成功率與穩定性。

\section{研究目的}

本研究旨在以 APIKGAgent 為基礎，設計可重複利用的成功經驗注入流程與局部修復流程，並分析這些改良對任務執行結果的影響。具體目標包括：萃取可遷移的成功經驗、建立動態修復策略、以及觀察經驗累積對後續任務的幫助。

\section{論文架構}

本論文共分為六章。第一章說明研究背景、動機與目的；第二章整理相關技術與文獻；第三章介紹系統方法與整體架構；第四章說明系統實作細節；第五章呈現實驗設計與結果分析；第六章總結研究成果並提出未來研究方向。
% 章節二 Chapter 2
\chapter{相關技術與文獻探討}

\section{大型語言模型代理人}

本節整理大型語言模型作為代理人核心時的基本概念、能力邊界與常見設計方式。

\section{API任務自動化}

本節整理 API 任務規劃、參數綁定、工具呼叫與多步驟執行流程的相關研究。

\section{案例式學習}

本節說明案例式學習如何透過既有成功案例支援推論、規劃與錯誤修正。

\section{動態計畫修復}

本節整理任務執行過程中，如何根據錯誤訊號與局部狀態修正原始計畫的相關方法。
% 章節三 Chapter 3
\begin{ZhChapter}

\chapter{CBR-APIKGAgent 系統架構與方法}

本章說明本研究所提出之 CBR-APIKGAgent 的系統架構與核心方法。承接前一章對多步驟 API 任務、經驗重用與動態計畫修復相關研究之整理，本章將重點放在兩項改良機制：其一為由訓練資料成功任務軌跡建立的原子化經驗重用機制，其二為針對執行失敗進行局部計畫調整的動態局部修復機制。以下首先介紹整體系統架構與任務流程，接著分別說明原子化經驗的表示、檢索與注入方式，以及局部修復模組如何根據失敗資訊進行插入、替換或刪除等修復操作。

\section{CBR-APIKGAgent 系統總覽}

CBR-APIKGAgent 的整體架構可視為在 APIKGAgent 原有流程中加入兩個輔助層：其一為離線原子化經驗庫與語意檢索流程，其二為執行失敗後的局部修復流程。以下先以整體架構圖呈現任務從初始規劃、步驟執行、局部修復到重新規劃的資料流與控制流，再整理主要元件之功能與相互關係。

\subsection{整體架構與主要元件}

CBR-APIKGAgent 以 APIKGAgent 為基礎，延續其 API 知識圖譜、API 搜尋、任務規劃、步驟式執行與重新規劃流程，並進一步加入原子化經驗重用與動態局部修復機制。整體而言，本研究並未改變 APIKGAgent 以「先搜尋 API、再產生計畫、最後逐步執行」為核心的任務執行模式，而是在規劃與錯誤處理階段加入案例式經驗檢索，使代理人在面對相似任務結構或局部執行失敗時，能取得來自訓練資料成功任務軌跡的策略參考。

CBR-APIKGAgent 的整體架構如圖 \ref{fig:cbr_apikgagent_architecture} 所示。系統以使用者任務為輸入，先透過 API 知識圖譜搜尋候選 API，並由任務規劃器產生初始任務計畫。任務執行期間，任務執行器會逐步執行計畫並保存中間結果；當步驟失敗時，系統會先嘗試由局部修復模組修正局部計畫，若局部修復無法處理，則回退至重新規劃器產生後續計畫。原子化經驗檢索器則同時支援初始規劃、重新規劃與局部修復三個決策階段，使代理人能取得與當前任務或錯誤情境相關的策略參考。

\begin{figure}[htbp]
    \centering
    \IfFileExists{figures/chapter3/cbr-apikgagent-architecture.png}{%
        \includegraphics[
            width=\textwidth,
            height=0.78\textheight,
            keepaspectratio
        ]{figures/chapter3/cbr-apikgagent-architecture.png}%
    }{%
        \fbox{\begin{minipage}{0.92\textwidth}
            \centering
            \textbf{CBR-APIKGAgent 整體架構示意圖}\\[0.8em]
            \begin{tabular}{p{0.24\textwidth} p{0.62\textwidth}}
                \hline
                \textbf{離線經驗層} &
                訓練資料成功任務軌跡
                $\rightarrow$
                原子化經驗庫
                $\rightarrow$
                規劃導向檢索／修復導向檢索 \\
                \hline
                \textbf{任務主流程} &
                使用者任務
                $\rightarrow$
                API 知識圖譜搜尋
                $\rightarrow$
                任務規劃器產生初始計畫
                $\rightarrow$
                任務執行器逐步執行 \\
                \hline
                \textbf{成功路徑} &
                步驟成功
                $\rightarrow$
                保存中間結果
                $\rightarrow$
                執行下一步
                $\rightarrow$
                完成任務 \\
                \hline
                \textbf{失敗處理} &
                步驟失敗
                $\rightarrow$
                形成錯誤證據
                $\rightarrow$
                局部修復模組（插入／替換／刪除）
                $\rightarrow$
                修復後回到任務執行器 \\
                \hline
                \textbf{回退機制} &
                局部修復失敗或超出範圍
                $\rightarrow$
                重新規劃器產生後續計畫
                $\rightarrow$
                回到任務執行器 \\
                \hline
            \end{tabular}\\[0.8em]
            \hrule
            \vspace{0.6em}
            圖中重點：原子化經驗檢索支援初始規劃、局部修復與重新規劃；局部修復優先處理可局部調整之錯誤，必要時再回退至重新規劃。
        \end{minipage}}%
    }
    \caption{CBR-APIKGAgent 整體架構}
    \label{fig:cbr_apikgagent_architecture}
\end{figure}

圖 \ref{fig:cbr_apikgagent_architecture} 中，案例式經驗檢索與動態局部修復雖然在任務失敗時會互相配合，但兩者在系統中扮演不同角色。案例式經驗檢索負責提供與當前任務或錯誤情境相關的策略參考，屬於決策輔助機制；動態局部修復則負責根據失敗資訊修改局部計畫，屬於錯誤處理機制。

CBR-APIKGAgent 的主要元件可分為六個部分：API 知識圖譜與 API 搜尋模組、任務規劃器、任務執行器、原子化經驗檢索器、動態局部修復模組，以及重新規劃器。為補充圖中各模組之功能，表 \ref{tab:cbr_apikgagent_components} 整理各元件在系統中的責任。

\begin{table}[htbp]
    \centering
    \caption{CBR-APIKGAgent 主要元件與功能}
    \label{tab:cbr_apikgagent_components}
    \small
    \renewcommand{\arraystretch}{1.15}
    \begin{tabularx}{0.95\textwidth}{
        >{\raggedright\arraybackslash}p{0.23\textwidth}
        >{\raggedright\arraybackslash}p{0.16\textwidth}
        >{\raggedright\arraybackslash}X
    }
        \hline
        \textbf{元件} & \textbf{角色} & \textbf{功能說明} \\
        \hline
        API 知識圖譜／搜尋模組 &
        API 資訊來源 &
        依據任務目標或修復需求搜尋相關 API，提供規劃與修復所需資訊。 \\[0.35em]
        任務規劃器 &
        初始規劃 &
        根據使用者任務、候選 API 與規劃經驗產生初始任務計畫。 \\[0.35em]
        任務執行器 &
        步驟執行 &
        依照任務計畫逐步執行 API 呼叫與資料處理，並保存中間結果與失敗資訊。 \\[0.35em]
        原子化經驗檢索器 &
        策略參考 &
        從離線經驗庫檢索相關經驗，提供初始規劃、重新規劃與局部修復作為策略參考。 \\[0.35em]
        局部修復模組 &
        錯誤修復 &
        在步驟失敗時，根據錯誤證據與修復經驗，以插入、替換或刪除步驟調整局部計畫。 \\[0.35em]
        重新規劃器 &
        回退處理 &
        當局部修復無法處理錯誤時，根據任務目標、既有計畫與失敗資訊重新產生後續計畫。 \\
        \hline
    \end{tabularx}
\end{table}

在資料流設計上，CBR-APIKGAgent 以任務狀態作為各階段之間共享上下文的載體。任務狀態保存任務目標、候選 API、目前計畫、已完成步驟結果與失敗資訊，使規劃、執行、局部修復與重新規劃能在同一脈絡下進行。較細部的狀態欄位、資料保存方式與程式傳遞流程，則於第四章系統實作中說明。

\subsection{架構擴充重點}

本研究以 APIKGAgent 作為基礎代理人架構，延續其以 API 知識圖譜輔助 API 搜尋、任務規劃、步驟式執行與重新規劃的基本流程。在此基礎上，CBR-APIKGAgent 並未改變原有的任務執行主軸，而是針對多步驟 API 任務中經驗重用不足與局部錯誤不易修復兩項問題進行擴充。

第一項擴充為原子化經驗重用機制。系統將訓練資料中的成功任務軌跡轉換為可檢索的經驗單元，並於初始規劃、重新規劃與局部修復階段提供與當前任務或錯誤情境相關的策略參考。此設計的目的不是直接複製過去任務軌跡，而是使代理人在面對相似資料依賴、規劃條件或常見陷阱時，能取得可供決策參考的經驗知識。

第二項擴充為動態局部修復機制。當任務執行過程中發生步驟失敗時，系統不立即進入整體重新規劃，而是先根據錯誤證據、相依資料、候選修復 API 與修復經驗評估是否存在可行的局部修復方式。若系統能產生局部調整方案，則嘗試透過插入、替換或刪除步驟修正局部計畫；若無法形成可行修復方案、修復次數達到限制，或錯誤訊號不足以支持局部修復，則回退至重新規劃。

此外，CBR-APIKGAgent 正式使用的經驗庫由訓練資料中的成功任務軌跡離線建立，測試期間不將新產生的修復紀錄寫回正式檢索索引。因此，本研究方法屬於離線經驗庫輔助的案例式修復架構，而非測試期間更新正式經驗庫的方法。

\subsection{任務執行與錯誤修復流程}

CBR-APIKGAgent 的任務執行流程可分為初始規劃、逐步執行、局部修復與重新規劃四個階段。各階段之功能如下：

\begin{enumerate}
    \item \textbf{初始規劃：}系統接收使用者任務後，透過 API 知識圖譜搜尋與任務相關的候選 API，並從原子化經驗庫中檢索規劃相關經驗。任務規劃器根據任務目標、候選 API 與檢索經驗產生初始任務計畫。

    \item \textbf{逐步執行：}產生初始計畫後，任務執行器依序執行計畫中的 API 呼叫與資料處理步驟。成功步驟的輸出會保存於任務狀態中，作為後續步驟的相依資料；若所有步驟皆成功完成，系統即依任務需求提交最終答案或完成環境狀態變更。

    \item \textbf{局部修復：}若某一步驟執行失敗，系統會收集失敗步驟、錯誤訊息、相關中間結果與 API 執行回饋，形成局部修復所需的錯誤證據。局部修復模組會結合錯誤證據、候選修復 API 與修復相關經驗，嘗試產生可行的局部修復方案。若修復成功，系統更新局部計畫並回到任務執行器繼續執行。

    \item \textbf{重新規劃：}若局部修復模組無法選出適合的修復 API、無法產生可執行的局部修復方案，或局部修復次數達到限制，系統則回退至重新規劃器。重新規劃器會參考原始任務、既有計畫、失敗資訊、已完成步驟結果與檢索經驗，重新產生後續計畫，並將修正後的計畫交回任務執行器繼續執行。
\end{enumerate}

整體而言，CBR-APIKGAgent 採取「先局部修復、必要時重新規劃」的錯誤處理流程。此流程的設計目的並非取代重新規劃，而是讓系統在錯誤僅侷限於局部步驟時，能先嘗試保留原本仍有效的計畫與資料流程；當局部修復不足以處理錯誤時，再回退至重新規劃。透過原子化經驗檢索與動態局部修復的結合，本研究進一步探討案例式策略參考是否有助於改善多步驟 API 任務中的規劃品質與錯誤修復表現。

\section{原子化經驗重用機制}

原子化經驗重用機制的目的，是將訓練資料中成功任務軌跡所隱含的規劃與資料流策略，轉換為可於不同任務間重用的經驗知識。相較於直接保存完整任務軌跡，本研究將經驗拆解為較小的策略單元，使其能作為初始規劃、重新規劃與局部修復階段的參考，而非要求代理人複製過去任務的完整步驟。

\subsection{經驗來源與使用範圍}

本研究使用的正式經驗庫由訓練資料中的成功任務軌跡離線建立，測試期間不新增經驗至正式經驗庫，也不更新正式檢索索引。此設計使經驗來源與實驗評估資料保持切割，避免測試期間產生的執行紀錄影響後續任務評估。

雖然系統在局部修復過程中可能產生修復紀錄，但本研究並未將此類紀錄納入正式經驗庫。主要原因在於，單一次局部修復紀錄不一定代表整體任務成功，且該修復是否正確需依後續資料流與最終任務結果判斷。若未經驗證即將修復紀錄轉換為可重用經驗，可能使錯誤或不完整的修復策略影響後續規劃。因此，本研究將測試期間的修復紀錄保留為分析資料，而非正式檢索經驗的一部分。

在經驗萃取階段，本研究以訓練資料中已知成功的任務為來源，參考任務規格、公開資料、標準解答、評估條件與實際 API 呼叫軌跡等資訊，萃取其中可跨任務重用的規劃原則與資料依賴模式。由於本研究關注的是多步驟 API 任務中的規劃與修復能力，因此經驗萃取時不以保留完整 API 呼叫序列為目標，而是著重於任務成功過程中可泛化的策略知識。

\subsection{原子化經驗定義}

本研究將原子化經驗定義為一個可跨任務重用的最小策略單元，用以描述特定任務情境下應保留的規劃原則、資料流模式、必要檢查與常見陷阱。換言之，案例式學習提供本研究的經驗重用觀點，而原子化經驗則是本研究用來表示與檢索案例知識的具體設計。

為避免與完整案例、API 模板或可執行技能混淆，原子化經驗在本研究中的定位如下：

\begin{itemize}
    \item \textbf{不同於完整案例}：原子化經驗不重播完整任務軌跡，而是抽取其中可跨任務重用的策略片段。
    \item \textbf{不同於 API 模板}：原子化經驗不保存固定 API 呼叫序列或參數值，而是描述資料流、檢查條件與決策原則。
    \item \textbf{不同於技能模組}：原子化經驗不是可直接執行的技能（skill），而是提供給規劃與修復流程參考的策略知識。
    \item \textbf{不同於線上長期記憶}：原子化經驗庫由訓練資料中的成功任務軌跡離線建立，測試期間不將新紀錄寫回正式經驗庫。
\end{itemize}

因此，每筆原子化經驗主要描述下列內容：

\begin{itemize}
    \item \textbf{適用情境}：何種任務或錯誤情境下可考慮使用此經驗。
    \item \textbf{核心原則}：當前任務應保留的主要規劃原則。
    \item \textbf{資料依賴}：需要維持的資料流、候選集合或目標集合關係。
    \item \textbf{常見陷阱}：應避免的規劃或修復方式。
\end{itemize}

具體而言，一筆原子化經驗通常對應到一類常見規劃需求或錯誤風險，例如完整取得多頁資料、從原始資料集合中建立正確目標集合、驗證候選結果是否符合任務來源、在寫入前比較目前狀態與目標狀態，或維持輸出格式契約等。這些經驗並不假設當前任務與來源任務完全相同，而是在任務結構、資料依賴或錯誤情境相似時，為代理人提供可參考的規劃與修復原則。

例如，若訓練任務要求代理人對某一條件下的多筆資料進行更新，原子化經驗不會保存「呼叫某個特定 API 並傳入某個固定參數」這類具體操作，而是萃取出「在執行寫入動作前，必須先由來源集合建立符合條件的目標集合，並確認每個目標項目的目前狀態」這類可跨任務重用的規劃原則。

綜合上述定義，原子化經驗在本研究中具有三項特性：

\begin{itemize}
    \item \textbf{粒度較小}：經驗粒度應小於完整任務軌跡，避免將來源任務的特定步驟或 API 名稱直接搬移到新任務。
    \item \textbf{內容可泛化}：經驗內容應描述可泛化的資料流與決策原則，例如目標集合建構、候選驗證、批次完整性與寫入保護。
    % TODO-FINAL-VERSION: 下列「適用性評估、規劃約束」偏向 v2／修正版經驗選擇設計。若最後採 v4，可改為較中性的「經檢索後作為策略參考或修復引導」；若採 v2，則保留並於第四章補上評分與約束注入細節。
    \item \textbf{作為策略約束與修復引導}：原子化經驗經適用性評估後，可轉換為規劃約束或修復引導，用以輔助計畫內容的檢查與調整。實際計畫與修復方案仍需依據當前任務、候選 API、錯誤資訊與執行狀態產生。
\end{itemize}

\subsection{原子化經驗表示方式}

為了使經驗能被檢索、排序並注入不同決策階段，本研究將每筆原子化經驗表示為結構化知識。表 \ref{tab:atomic_experience_schema} 中的欄位用以說明本研究如何將成功任務軌跡抽象化為可檢索經驗；實際資料格式、儲存方式與載入流程則於第四章系統實作中說明。

\begin{table}[htbp]
    \centering
    \caption{原子化經驗主要欄位與用途}
    \label{tab:atomic_experience_schema}
    \small
    \renewcommand{\arraystretch}{1.15}
    \begin{tabularx}{0.95\textwidth}{
        >{\raggedright\arraybackslash}p{0.25\textwidth}
        >{\raggedright\arraybackslash}X
    }
        \hline
        \textbf{欄位} & \textbf{用途說明} \\
        \hline
        \texttt{experience\_type} &
        經驗類型標籤，用於描述此經驗所屬的粗略策略類別，例如目標集合建構、候選驗證、批次完整性或輸出格式契約。 \\[0.35em]
        \texttt{source\_task\_id} &
        來源訓練任務識別碼，用於追溯經驗來源與後續分析。 \\[0.35em]
        \texttt{source\_task\_goal} &
        來源訓練任務目標，用於語意檢索與理解經驗適用背景。 \\[0.35em]
        \texttt{when\_to\_use} &
        描述此經驗適用的任務情境或錯誤情境。 \\[0.35em]
        \texttt{core\_principle} &
        此經驗欲保留的核心規劃原則或資料流不變條件。 \\[0.35em]
        \texttt{key\_patterns} &
        此經驗所包含的關鍵資料流或規劃模式，用於提醒代理人應保留的中間集合、驗證條件或操作順序。 \\[0.35em]
        \texttt{how\_to\_apply} &
        % TODO-FINAL-VERSION: how_to_apply 已於 v4 補入，但其是否用於索引、直接提示注入，或 v2 的適用性評分／約束化流程，需依最終版本於第四章寫清楚。
        將此經驗套用到新任務時可參考的抽象步驟，但不對應到特定 API 呼叫序列。 \\[0.35em]
        \texttt{pitfalls} &
        此經驗欲避免的常見錯誤，例如過早執行寫入、忽略篩選條件、未完整取得資料或使用錯誤目標集合。 \\
        \hline
    \end{tabularx}
\end{table}

上述欄位使原子化經驗同時具備可追溯性、可檢索性與可注入性。其中，來源任務資訊主要用於建立檢索語境；適用情境、核心原則、關鍵模式與常見陷阱則用於幫助代理人在新任務中辨識相似規劃結構。為避免經驗過度貼近單一訓練任務，本研究在經驗表示上避免保留具體 API 名稱、固定參數值或完整任務流程，而以較抽象的資料來源、候選集合、驗證條件、目標集合與輸出契約等概念描述經驗。

\subsection{經驗檢索流程}

CBR-APIKGAgent 透過語意檢索機制，在不同任務階段取得與當前情境相關的原子化經驗。由於初始規劃與錯誤修復所面對的資訊不同，本研究將同一份離線經驗庫建立為兩種檢索視角：規劃導向索引與修復導向索引。

兩種檢索視角的用途如下：

\begin{itemize}
    \item \textbf{規劃導向索引：}用於初始規劃與重新規劃階段，主要以經驗中的來源任務目標、適用情境、核心原則、關鍵模式與套用方式建立比對內容。此索引的目標，是協助系統根據目前任務目標或既有計畫，找到在資料依賴、目標集合建構或操作順序上相似的經驗。
    \item \textbf{修復導向索引：}用於局部修復與重新規劃中的錯誤分析，主要以經驗中的適用情境、核心原則、關鍵模式、套用方式與常見陷阱建立比對內容。此索引的目標，是使系統能根據失敗訊息或錯誤證據，取得與當前錯誤情境相關的修復參考。
\end{itemize}

換言之，檢索時的語意比對包含兩個部分：一方面，系統會依目前所處階段產生查詢內容，例如任務目標、失敗訊息或重新規劃脈絡；另一方面，經驗庫中的每筆原子化經驗會依規劃導向或修復導向索引，選取不同欄位形成可比對的經驗表示。透過此設計，同一筆經驗可依使用情境呈現不同的檢索重點。

依照使用階段不同，經驗檢索流程可分為以下三種情境：

\begin{enumerate}
    \item \textbf{初始規劃階段}

    \begin{itemize}
        \item \textbf{查詢內容：}以任務目標作為主要查詢。
        \item \textbf{檢索視角：}使用規劃導向索引取得候選原子化經驗。
        % TODO-FINAL-VERSION: 「適用性評估」是偏 v2／修正版的概念。若最後採 v4，可弱化為「依語意相關性與任務階段選取經驗」；若採 v2，保留並在第四章補 LLM applicability scoring。
        \item \textbf{適用性評估：}判斷候選經驗是否能對當前任務形成具體且安全的規劃約束，例如資料依賴、目標範圍、候選驗證、操作順序、批次完整性或輸出限制等。
        % TODO-FINAL-VERSION: 「通過評估、轉換為規劃約束、檢查計畫是否反映約束」偏 v2。若最後採 v4，改成「檢索經驗會被整理並注入提示，作為規劃時的策略參考」。
        \item \textbf{使用方式：}通過評估的經驗會被轉換為規劃約束，注入初始規劃提示中，並用於輔助檢查產生之計畫是否反映相關約束。
    \end{itemize}

    \item \textbf{局部修復階段}

    \begin{itemize}
        \item \textbf{查詢內容：}以失敗訊息、失敗步驟與當前錯誤脈絡作為查詢。
        \item \textbf{檢索視角：}使用修復導向索引取得與錯誤情境相關的候選經驗。
        % TODO-FINAL-VERSION: 「適用性評估」偏 v2／修正版。若最後採 v4，可改為「依錯誤情境檢索並提供相關修復經驗」。
        \item \textbf{適用性評估：}判斷經驗是否有助於解釋當前失敗或形成具體修復依據，例如補足前置資料、修正操作目標、維持目標來源、補齊部分執行、避免副作用或處理其他與錯誤脈絡相關的修復需求。
        % TODO-FINAL-VERSION: 「通過評估」偏 v2；v4 可保留「檢索經驗被整理為修復引導」但不強調評分門檻。
        \item \textbf{使用方式：}通過評估的經驗會被整理為修復引導，供局部修復流程產生插入、替換或刪除步驟時參考。
    \end{itemize}

    \item \textbf{重新規劃階段}

    \begin{itemize}
        \item \textbf{查詢內容：}同時考量任務目標、既有計畫、失敗步驟與錯誤原因。
        \item \textbf{檢索視角：}同時使用規劃導向與修復導向索引，分別取得與任務結構及失敗情境相關的候選經驗。
        % TODO-FINAL-VERSION: 重新規劃階段的「分別評估」偏 v2／修正版。若最後採 v4，改為「分別取得並整理」即可。
        \item \textbf{適用性評估：}分別評估規劃經驗與修復經驗是否適用於當前重新規劃情境。
        % TODO-FINAL-VERSION: 「規劃約束與修復引導」可支援 v2；若 v4 final，建議弱化為「策略參考與修復提醒」。
        \item \textbf{使用方式：}通過評估的經驗會被整合為重新規劃時的規劃約束與修復引導，使新計畫能同時回應原始任務需求與已觀察到的錯誤原因。
    \end{itemize}
\end{enumerate}

% TODO-FINAL-VERSION: 本段最後一句提到「適用性評估具體評分方式、候選經驗篩選門檻、約束檢查」，偏 v2／修正版。若 final 採 v4，需改為 direct retrieval + prompt injection + how_to_apply；若採 v2，第四章補 scoring、selected experiences、constraints 與檢查流程。
整體而言，此檢索流程的重點在於讓不同階段以不同角度使用同一批離線經驗。初始規劃偏重「任務結構是否相似」，局部修復偏重「錯誤情境是否相似」，重新規劃則同時參考規劃導向與修復導向經驗，以協助代理人在修正後續計畫時兼顧原任務目標與已觀察到的失敗原因。原子化經驗重用機制提供的是可檢索並可轉換為規劃約束或修復引導的案例式知識，而非固定規則或保證正確的修復模板；適用性評估的具體評分方式、候選經驗篩選門檻、提示注入位置，以及計畫產生後如何檢查經驗約束是否被反映，則於第四章系統實作中進一步說明。

\section{動態局部修復機制}

動態局部修復機制的目的，是在任務執行過程中發生局部失敗時，先嘗試針對失敗步驟附近的計畫進行小範圍調整，而不是立即重新產生整體計畫。此設計的核心假設是：部分 API 任務失敗並非源自整體任務理解錯誤，而可能是缺少前置資料、使用不適合的 API、資料依賴未正確銜接，或某一步驟的操作目標不完整所造成。若此類問題能透過局部插入、替換或刪除步驟處理，便有機會保留原計畫中已正確完成的部分，並減少重新規劃對既有資料流程造成干擾的可能性。

在本研究所探討的多步驟 API 任務中，代理人執行過程可能已對環境狀態造成改變。因此，錯誤處理流程並不假設系統能將環境回復至任務初始狀態或錯誤發生前狀態。無論是既有重新規劃機制或本研究提出的局部修復機制，皆是在目前已觀察到的執行狀態下產生後續行動；兩者差異在於，重新規劃傾向重新產生較大範圍的後續計畫，而局部修復則優先嘗試對失敗步驟附近進行小範圍調整。

基於上述設定，本研究不以可驗證的錯誤根因診斷作為局部修復之前提。若要精確驗證不同錯誤根因假設，通常需要能夠回復環境狀態、重放特定步驟，或比較多個修復分支；然而，此類流程較接近離線除錯或搜尋式求解，並非本研究所關注的單次任務執行脈絡下之代理人錯誤恢復能力。需要注意的是，局部修復並不保證每次皆能改善任務結果。本研究將局部修復視為一種受限制的錯誤恢復嘗試：系統根據目前可取得的錯誤訊息、失敗步驟、相依資料與案例式修復經驗，判斷是否存在可行的局部調整方案；若證據不足或修復嘗試超出可控範圍，則回退至重新規劃流程。

\subsection{錯誤資訊與相依資料蒐集}

局部修復的品質高度依賴失敗發生時可取得的上下文資訊。若系統只根據單一錯誤訊息選擇修復 API，容易忽略原任務目標、前後步驟關係或已完成資料流程，導致修復方向偏離原計畫。因此，本研究在步驟失敗後，會先於任務執行器端整理錯誤訊息並產生暫定錯誤分類；若錯誤處理路由判斷可進入局部修復，修復流程再進一步使用這些資訊搜尋候選 API、檢索案例式修復參考並產生局部計畫修改。

依照實際處理順序，錯誤資訊與相依資料蒐集可分為下列四個步驟：

\begin{enumerate}
    \item \textbf{記錄失敗步驟與原始錯誤訊息：}當某一步驟執行失敗時，系統會保留該步驟在整體計畫中的位置、步驟描述、原本欲完成的局部目標，以及 API 呼叫或資料處理所回傳的錯誤訊息。
    \item \textbf{整理 API 回饋與相依資料脈絡：}系統會進一步保存失敗 API 的相關資訊、錯誤細節、前序步驟產生的中間結果，以及目前可觀察到的資料依賴脈絡。這些資訊用於判斷失敗是否可能與缺少前置資料、資料來源不一致或操作目標不完整有關。
    \item \textbf{產生暫定錯誤分類與可行回饋：}在錯誤處理路由做出決策前，系統會參考 AgentErrorTaxonomy \cite{zhu2025agentdebug} 所整理之代理人錯誤分類概念，並根據失敗訊息、失敗步驟、API 回饋與相依資料，由大型語言模型產生兩項核心輸出：

    \begin{itemize}
        \item \textbf{暫定錯誤分類（\texttt{error\_type}）：}描述目前失敗較可能接近何種錯誤情境，例如參數缺漏、資料查詢失敗、權限或狀態不符，或原步驟可能需要拆分與較大範圍調整等。在本研究實作中，若原始失敗步驟被暫定分類為不可行計畫，錯誤處理路由會直接回退至重新規劃。
        \item \textbf{可行回饋（\texttt{actionable\_feedback}）：}將錯誤訊息整理為後續修復可使用的操作建議，例如補足前置資料、重新搜尋候選 API，或避免重複執行已完成動作。此回饋不直接決定是否進入局部修復；若系統進入局部修復流程，則會作為修復 API 搜尋與修復策略選擇的重要依據。
    \end{itemize}

    \item \textbf{交由錯誤處理路由與局部修復流程使用：}錯誤處理路由會先根據暫定錯誤分類與修復限制判斷是否進入局部修復。若進入局部修復，上述錯誤證據與可行回饋會被用於修復 API 搜尋；後續修復策略選擇時，系統再結合候選 API、相依資料與案例式修復參考，判斷應插入、替換或刪除局部計畫步驟。
\end{enumerate}

需要說明的是，本研究雖參考錯誤分類研究整理失敗資訊，但此錯誤情境判斷並不等同於可驗證的錯誤根因診斷，也不代表系統能正確辨識真正的上游錯誤原因。多步驟 API 任務中的錯誤原因往往需要結合失敗訊息、原步驟目的、前序資料與候選 API 才能推估；同一個執行錯誤也可能來自不同上游原因。因此，本研究將錯誤情境判斷視為修復流程中的輔助訊號，而非決定修復結果的唯一依據。後續修復策略選擇時，系統會再次結合候選 API、案例式修復經驗與相依資料進行判斷，而不是只依據初步錯誤分類直接決定修復方式。

\subsection{錯誤處理路由與局部修復觸發條件}

在任務執行器完成錯誤資訊整理後，系統會透過錯誤處理路由判斷是否進入局部修復流程。因此，局部修復並非由多種獨立事件任意觸發，而是以「步驟執行失敗」作為主要觸發點，再根據暫定錯誤分類、失敗步驟類型、修復嘗試次數與修復鏈深度決定後續方向。

在本研究架構中，錯誤處理路由會先判斷目前失敗步驟是否屬於修復流程新增或替換出的步驟。之所以需要區分這兩類情境，是因為修復步驟再次失敗時，不應被視為新的獨立失敗點，而應回掛至其原本欲修復的原始失敗步驟，並沿用同一失敗點的修復嘗試次數與修復鏈深度限制。依此路由規則，進入局部修復的情境可整理為兩類：

\begin{itemize}
    \item \textbf{原始計畫步驟失敗：}當原始計畫中的 API 呼叫或資料處理步驟回傳錯誤，或其結果不足以支撐後續計畫時，若該步驟未被暫定分類為不可行計畫，且該步驟尚未達到局部修復嘗試上限，系統會優先嘗試局部修復。
    \item \textbf{修復步驟再次失敗：}若先前插入或替換出的修復步驟執行失敗，系統會透過修復步驟與原始步驟之間的對應關係，回到其欲修復的原始失敗點，檢查該原始失敗點是否仍有修復嘗試額度與修復鏈深度空間；若尚未超出限制，則再次嘗試產生新的局部修復方案。
\end{itemize}

上述兩類情境皆以「同一失敗點仍在可控修復範圍內」為前提。若原始失敗步驟被暫定分類為不可行計畫，或同一原始失敗點已達修復嘗試上限與修復鏈深度限制，系統即不再進入局部修復，而是回退至重新規劃。換言之，錯誤處理路由主要依據失敗步驟是否為修復步驟、暫定錯誤分類中的不可行計畫訊號，以及局部修復嘗試限制，決定後續應進入局部修復或重新規劃；可行回饋不直接決定路由，而是在進入局部修復後作為修復行動的依據。

\subsection{修復 API 搜尋}

在取得錯誤證據與可行回饋後，系統會重新搜尋可能用於修復的候選 API。與初始規劃階段以使用者任務為主要查詢不同，修復 API 搜尋的查詢重點更接近「如何讓目前失敗步驟恢復可執行狀態」或「是否有其他 API 能完成失敗步驟原本應達成的局部目的」。

本研究將修復階段使用的資訊分為「搜尋查詢」與「後續判斷上下文」兩個層次：

\begin{itemize}
    \item \textbf{搜尋查詢：}主要由失敗訊息與可行回饋組成。失敗訊息提供目前執行失敗的表面徵象，可行回饋則將錯誤資訊轉換為較具操作方向的修復需求，例如補足缺漏資料、重新取得目標物件，或避免不適當的 API 呼叫。
    \item \textbf{後續判斷上下文：}包含失敗步驟內容與其局部目標、候選 API 說明、結構化失敗證據，以及修復導向經驗。其中，結構化失敗證據可包含錯誤代碼、錯誤細節、失敗 API 名稱、結果分析、相依資料脈絡與 API 呼叫經驗等資訊。這些資訊主要用於後續修復策略選擇，而不是單純全部併入 API 搜尋查詢中。
\end{itemize}

透過上述區分，搜尋階段可聚焦於取得可能有用的工具；在取得候選 API 後，系統會進一步檢索修復導向經驗，並將候選 API、結構化失敗資訊與修復經驗一併提供給局部修復模組，用以判斷候選 API 是否適合插入、替換或刪除原計畫步驟。

修復 API 搜尋並不代表系統必然會加入新的 API。候選 API 取得後，局部修復模組仍需在同一個策略選擇過程中，根據錯誤證據、失敗步驟內容、候選 API、相依資料脈絡與修復經驗，判斷候選 API 與原失敗步驟之間的關係。舉例而言，候選 API 可能只是用來補足原步驟缺少的前置資料，也可能直接承接原步驟欲完成的操作；在部分情境下，錯誤證據亦可能顯示原步驟本身不應繼續執行。此中間判斷的目的，是降低後續選擇插入、替換或刪除策略時的模糊性，而非建立完整的修復意圖分類。

\subsection{局部計畫修復策略}

本研究採用插入步驟（Insert）、替換步驟（Replace）與刪除步驟（Delete）三種局部計畫修復策略。三種策略皆以最小化計畫變動為原則，僅調整失敗步驟附近的計畫片段，並盡可能保留已成功執行的步驟與可用中間資料。實際策略選擇時，系統會將錯誤證據、候選修復 API、結構化失敗資訊與修復導向經驗提供給大型語言模型，由其判斷修復步驟與原失敗步驟之間的關係，例如是否應保留原失敗步驟繼續執行、是否應以新 API 取代原步驟，或是否應移除原步驟。

\begin{enumerate}
    \item \textbf{插入步驟（Insert）}

    當失敗原因可能來自缺少前置資料或必要條件時，系統會考慮在原失敗步驟前插入新的修復步驟。此策略適用於原步驟本身仍可能正確，但執行前缺少必要輸入的情境。例如，原步驟需要某個物件識別碼，但計畫尚未先查詢符合條件的物件；此時修復模組可插入查詢步驟，再讓原失敗步驟重新執行。

    插入策略的重點在於保留原失敗步驟的角色。新增步驟只負責補足資料或建立前置條件，而不應任意改變原步驟的任務語意。因此，若候選 API 與原步驟具有相同操作目的，而非僅提供前置資料，系統需進一步評估是否應改採替換策略。

    \item \textbf{替換步驟（Replace）}

    當錯誤訊號顯示原失敗步驟本身可能使用了不適當的 API、錯誤的操作方式，或其描述無法完成原本的局部目標時，系統會考慮以新的修復步驟替換原步驟。此策略適用於候選 API 能直接承接原步驟目的，且保留原步驟反而可能造成重複執行或錯誤副作用的情境。

    替換策略比插入策略更可能影響資料流程，因此本研究採取較保守的使用原則。系統需根據失敗證據、候選 API 功能與相依資料，確認新步驟能合理承接原步驟在計畫中的位置；若候選 API 僅能提供輔助資料，則較適合使用插入策略，而非直接取代原步驟。

    \item \textbf{刪除步驟（Delete）}

    當失敗步驟被判斷為冗餘、重複或不應繼續執行時，系統會考慮將其從局部計畫中移除。此策略適用於原步驟與目前任務狀態不一致，或該步驟可能重複處理已完成操作的情境。

    刪除策略的風險在於可能移除後續步驟仍需要的資料來源，因此系統在刪除步驟後，仍需維持計畫中步驟順序與相依關係的一致性。若刪除後導致後續計畫缺少必要資料，則局部修復不應被視為可行方案。
\end{enumerate}

無論採用何種策略，局部修復完成後，系統皆需更新局部計畫，使後續執行流程能接續修復後的步驟順序與資料依賴。此處的更新包含調整步驟位置、維持或重建相依關係，以及確認修復步驟與原任務目標仍保持一致。實際計畫修改、步驟索引調整與狀態追蹤方式，將於第四章系統實作中說明。

\subsection{局部修復與重新規劃的切換條件}

局部修復與重新規劃並非互相排斥的機制，而是錯誤處理流程中的兩個層次。局部修復負責處理可由小範圍計畫調整嘗試恢復的失敗；重新規劃則作為回退機制，用於處理局部修復無法形成可行方案或修復嘗試失敗的情境。

本研究將下列情況視為切換至重新規劃的依據：

\begin{itemize}
    \item \textbf{原始步驟被判斷為不可行計畫：}當原始失敗步驟被暫定分類為不可行計畫時，系統不嘗試局部修復，而是直接回退至重新規劃。
    \item \textbf{缺乏可用修復 API 或刪除依據：}若候選 API 無法補足前置條件、替代原步驟或完成缺漏操作，且失敗步驟也不適合直接刪除，局部修復流程無法產生計畫變更。
    \item \textbf{修復嘗試達到限制：}同一失敗位置經過有限次局部修復後仍無法恢復執行，為避免重複嘗試造成迴圈，系統停止局部修復。
    \item \textbf{修復鏈深度達到限制：}若修復步驟本身再次失敗，且同一失敗點已累積過多層修復步驟，系統停止繼續擴張局部修復鏈。
\end{itemize}

透過上述切換條件，CBR-APIKGAgent 採取「優先嘗試局部修復，但保留重新規劃回退」的錯誤處理策略。此設計一方面使系統能在局部錯誤發生時嘗試保留既有正確步驟，另一方面也避免在錯誤原因不明或修復方向不可靠時持續修改同一段計畫。第五章將進一步透過實驗與案例分析，評估此局部修復機制對任務完成能力與錯誤修復表現的影響。

\end{ZhChapter}

% 章節四 Chapter 4
\begin{ZhChapter}

\chapter{系統實作與細節}

\section{開發環境與工具}

在此描述實驗與開發時使用的環境、語言、框架與工具。

\section{原 APIKGAgent 基礎執行流程與本研究修改點}

\subsection{Main Graph / Executor Graph / API Execute Graph 概述}

\subsection{Process Step / API Step 概述}

\subsection{本研究延伸位置說明}

\section{全域成功經驗建立與注入}

\subsection{train dataset 成功經驗萃取}

\subsection{RAG 語意檢索與 Planner 注入}

\section{Rerollout 動態修復實作}

\subsection{錯誤分類與修復 API 搜尋}

\subsection{任務步驟的新增、刪除與更新}

\subsection{與原 replan 機制之差異}

\section{區域成功修復經驗儲存與檢索}

\subsection{成功修復經驗萃取}

\subsection{Local recovery lessons 儲存與後續任務再利用}

\section{其他實作細節}

在此補充其他實作中需要交代的細節。

\end{ZhChapter}

% 章節五 Chapter 5
\begin{ZhChapter}

\chapter{實驗與結果分析}

本章旨在評估 CBR-APIKGAgent 於多步驟 API 任務中的任務完成能力與錯誤修復表現。透過 AppWorld 任務環境進行實驗，觀察完整系統之整體表現，並進一步分析原子化經驗重用與動態局部修復機制各自對任務完成率造成的影響。

由於本研究的核心假設為「由訓練資料成功任務軌跡萃取之原子化經驗，可作為相似任務結構或錯誤情境下的策略參考」，以及「局部修復可在部分失敗情境下協助代理人由執行錯誤中恢復」，因此本章將分別從整體效能、消融實驗與案例分析等面向進行分析。實驗結果以實際執行紀錄與 AppWorld 評估結果為依據，並搭配可取得完整軌跡的案例，討論兩項機制如何在不同階段共同影響任務完成，以及完整系統目前呈現的限制。

\section{研究問題}

本研究設計四項研究問題（Research Questions, RQs），分別以 RQ1 至 RQ4 表示，以對應第一章所提出之研究目的：

\begin{enumerate}
    \item \textbf{RQ1：CBR-APIKGAgent 相較於既有方法在 AppWorld 任務中的整體表現為何？}

    此問題用於評估完整系統在 AppWorld \texttt{test\_normal} 上的整體任務完成能力。

    \item \textbf{RQ2：原子化經驗重用機制對任務完成率造成何種影響？}

    此問題用於透過消融實驗分析原子化經驗重用機制對任務完成結果的影響。

    \item \textbf{RQ3：動態局部修復機制對任務完成率造成何種影響？}

    此問題用於透過消融實驗分析動態局部修復機制對任務完成結果的影響。

    \item \textbf{RQ4：原子化經驗與動態局部修復如何共同影響多步驟 API 任務的完成？}

    此問題延續消融實驗結果，透過可取得完整軌跡的成功與限制案例，觀察原子化經驗與動態局部修復分別在規劃與執行階段的作用，以及兩者如何共同影響任務完成。
\end{enumerate}

\section{實驗設置}

\subsection{資料集與任務範圍}

本研究於 AppWorld \cite{trivedi2024appworld} 環境中進行實驗。AppWorld 提供多應用程式 API 任務環境，任務完成與否由最終環境狀態是否符合預期條件判定。因此，代理人不僅需要產生合理文字回覆，也必須透過 API 呼叫與資料處理正確改變或查詢環境狀態。

正式實驗以 AppWorld \texttt{test\_normal} 作為主要評估資料集，共包含 168 個任務，並以本研究提出之 CBR-APIKGAgent 作為評估系統。各項實驗所使用的任務範圍將於對應實驗小節中說明。

在每一項實驗設定中，各測試任務僅採用一次正式執行結果；局部修復、重新規劃與執行重試皆發生於該次任務執行期間，並受後述控制設定限制。任務結束後，AppWorld evaluation tests 的結果僅用於計算評估指標，不會回饋給代理人、寫入正式經驗庫，或作為調整後重新執行同一測試任務的依據。

\subsection{評估指標}

依據 AppWorld 官方評估方式，AppWorld 採用程式化且以最終狀態為基礎的 evaluation tests 判斷任務是否完成，並以任務目標完成率（Task Goal Completion, TGC）與情境目標完成率（Scenario Goal Completion, SGC）作為主要指標 \cite{trivedi2024appworld}。

\begin{enumerate}
    \item \textbf{任務目標完成率（TGC）：}TGC 表示通過所有 AppWorld evaluation tests 的任務比例。對於單一任務而言，只有當該任務的所有評估測試皆通過時，該任務才計為成功。其計算方式如式 \ref{eq:tgc} 所示：

    \begin{equation}
        \mathrm{TGC} = \frac{N_{\mathrm{success}}}{N_{\mathrm{total}}} \times 100\%
        \label{eq:tgc}
    \end{equation}

    其中，\(N_{\mathrm{success}}\) 代表通過所有 evaluation tests 的任務數量，\(N_{\mathrm{total}}\) 代表評估任務總數。

    \item \textbf{情境目標完成率（SGC）：}
    SGC 以 AppWorld 的 task scenario 為單位計算。只有同一 scenario 產生之所有任務皆通過 evaluation tests 時，該 scenario 才計為成功。此指標用於衡量代理人是否能在同一任務情境的不同需求與初始狀態變化下穩定完成目標。其計算方式如式 \ref{eq:sgc} 所示：

    \begin{equation}
        \mathrm{SGC} = \frac{N_{\mathrm{success\ scenarios}}}{N_{\mathrm{total\ scenarios}}} \times 100\%
        \label{eq:sgc}
    \end{equation}
\end{enumerate}

\subsection{實驗環境與主要控制設定}

本研究實驗之硬體與軟體環境如表 \ref{tab:hardware_software_environment} 所示。由於本研究主要透過雲端大型語言模型服務進行推論，本地端硬體主要負責執行 AppWorld 環境、代理人控制流程、API 模擬環境與實驗紀錄整理。

\begin{table}[htbp]
    \centering
    \caption{硬體與軟體環境}
    \label{tab:hardware_software_environment}
    \small
    \renewcommand{\arraystretch}{1.15}
    \begin{tabularx}{0.82\textwidth}{
        >{\raggedright\arraybackslash}p{0.3\textwidth}
        >{\raggedright\arraybackslash}X
    }
        \hline
        \textbf{項目} & \textbf{設定} \\
        \hline
        作業系統 & Ubuntu 22.04.5 LTS（WSL2） \\
        處理器 & 12th Gen Intel(R) Core(TM) i7-12700 \\
        記憶體 & 16 GB \\
        Python 版本 & Python 3.12 \\
        AppWorld 版本 & 0.1.3 \\
        \hline
    \end{tabularx}
\end{table}

除上述執行環境外，實驗亦需固定語言模型、推論溫度、執行預算、經驗檢索設定與局部修復限制等主要控制項。本研究採用 Gemini 2.5 Flash 作為大型語言模型 \cite{geminiteam2025gemini25}；表 \ref{tab:experiment_environment} 列出本研究實驗中需要固定或比較的重點設定。

\begin{table}[htbp]
    \centering
    \caption{實驗主要控制設定}
    \label{tab:experiment_environment}
    \small
    \renewcommand{\arraystretch}{1.15}
    \begin{tabularx}{0.82\textwidth}{
        >{\raggedright\arraybackslash}p{0.3\textwidth}
        >{\raggedright\arraybackslash}X
    }
        \hline
        \textbf{項目} & \textbf{設定} \\
        \hline
        大型語言模型 & \texttt{gemini-2.5-flash} \\
        溫度參數 & 判斷 0；規劃、重新規劃與執行 0.2；API 評估 0.1 \\
        任務執行時間限制 & 3600 秒／題 \\
        執行流程內重新嘗試上限 & 3 次／題 \\
        局部修復上限 & 3 次／原始失敗步驟 \\
        修復鏈深度上限 & 2 步驟／原始失敗步驟 \\
        經驗檢索相似度門檻 & 0.35 \\
        初始規劃經驗檢索 & 規劃導向至多 3 筆 \\
        重新規劃經驗檢索 & 規劃導向至多 1 筆；修復導向至多 1 筆 \\
        局部修復經驗檢索 & 修復導向至多 2 筆 \\
        \hline
    \end{tabularx}
\end{table}

\section{實驗一：整體效能比較}
\label{sec:overall_performance}

本節對應 RQ1，目的在於呈現 CBR-APIKGAgent 在 AppWorld \texttt{test\_normal} 上的整體任務完成率，並與既有方法進行比較。本實驗使用完整 \texttt{test\_normal} 168 個任務作為評估範圍，並依 AppWorld 官方評估結果計算 TGC 與 SGC。若任務執行後未產生完整評估報告或出現異常執行情形，該任務在統計中計為失敗。

為呈現本研究方法於 AppWorld 基準中的相對位置，表 \ref{tab:overall_performance} 列出既有方法於 \texttt{test\_normal} 上的結果。表中 GPT-4o 與 LLaMA3-70B 設定下之 ReAct、PlanExec、FullCodeRefl 與 IPFunCall 結果取自 AppWorld 原論文 \cite{trivedi2024appworld}；Qwen2.5-32B 設定下之 ReAct，以及 Gemini 2.5 Flash 設定下之 ReAct、FullCodeRefl 與 APIKGAgent 結果取自 APIKGAgent 原研究整理 \cite{wu2025apikgagent}；IBM CUGA 之結果取自其原研究報告 \cite{shlomov2026benchmarks}；LOOP 之結果取自本研究整理時可取得之 AppWorld 官方 leaderboard \cite{appworld_leaderboard}，其方法背景則參考 LOOP 原研究 \cite{chen2025loop}。其中，leaderboard 數據以 2026 年 7 月 13 日存取結果為準。CBR-APIKGAgent 為本研究於相同 \texttt{test\_normal} 任務集合下之實驗結果。

由於不同方法可能採用不同模型、示範資料、訓練策略或執行預算，因此表 \ref{tab:overall_performance} 主要作為 benchmark 相對表現參考，而非嚴格控制變因下的直接比較。各機制對本研究系統造成的影響，將於後續消融實驗中以相同模型與相同任務集合進一步分析。

\begin{table}[htbp]
    \centering
    \caption{各方法於 AppWorld \texttt{test\_normal} 之整體效能比較}
    \label{tab:overall_performance}
    \small
    \renewcommand{\arraystretch}{1.15}
    \begin{tabularx}{0.8\textwidth}{
        >{\raggedright\arraybackslash}X
        >{\raggedright\arraybackslash}p{0.22\textwidth}
        >{\centering\arraybackslash}p{0.12\textwidth}
        >{\centering\arraybackslash}p{0.12\textwidth}
    }
        \hline
        \textbf{Method} & \textbf{Model} & \textbf{TGC (\%)} & \textbf{SGC (\%)} \\
        \hline
        ReAct & \multirow{4}{*}{GPT-4o} & 48.8 & 32.1 \\
        PlanExec &  & 44.6 & 23.2 \\
        FullCodeRefl &  & 33.9 & 26.8 \\
        IPFunCall &  & 32.1 & 16.1 \\
        \hline
        IBM CUGA & GPT-4.1 & 73.2 & 62.5 \\
        \hline
        ReAct & \multirow{3}{*}{LLaMA3-70B} & 20.8 & 8.9 \\
        PlanExec &  & 8.9 & 1.8 \\
        FullCodeRefl &  & 24.4 & 17.9 \\
        \hline
        ReAct & \multirow{2}{*}{Qwen2.5-32B} & 39.2 & 18.6 \\
        LOOP &  & 72.6 & 53.6 \\
        \hline
        ReAct & \multirow{4}{*}{Gemini 2.5 Flash} & 50.6 & 28.6 \\
        FullCodeRefl &  & 36.9 & 26.8 \\
        APIKGAgent &  & 41.7 & 17.9 \\
        \textbf{CBR-APIKGAgent} &  & \textbf{51.8} & \textbf{35.7} \\
        \hline
    \end{tabularx}
\end{table}

由表 \ref{tab:overall_performance} 可知，CBR-APIKGAgent 於 \texttt{test\_normal} 上達成 51.8\% 的 TGC 與 35.7\% 的 SGC。在表中採用 Gemini 2.5 Flash 的方法裡，本研究方法之兩項指標皆為最高，並高於同模型的 ReAct、FullCodeRefl 與 APIKGAgent；但與 LOOP 及 IBM CUGA 相比仍有差距。

從方法設計來看，LOOP 透過強化學習直接在 AppWorld 環境中調整代理人的行為，使 API 文件查閱、行動選擇與錯誤恢復策略反映於模型的互動行為；本研究則不調整模型權重，而是在推論階段檢索原子化經驗，並將其加入規劃與修復脈絡以引導模型重用既有策略。因此，本研究方法的效果仍會受到經驗檢索結果、任務限制契合度及模型對經驗內容之理解所影響，此項學習方式的差異可能是 LOOP 取得較高完成率的原因之一。另一方面，CUGA 採用階層式規劃器與執行器架構，由持久化任務紀錄維持步驟、變數綁定與重新規劃狀態，再將子任務交由專用代理人處理；本研究則著重於在 APIKGAgent 流程中加入經驗重用與失敗後的局部修復。從本研究的限制案例可觀察到，當任務涉及多階段資料篩選或搜尋候選結果時，原始任務條件與後續操作目標之間的對應仍可能產生偏差，因此，跨步驟狀態維持方式的差異可能是 CUGA 取得較高完成率的原因之一。

然而，LOOP、CUGA 與本研究使用的基礎模型、學習方式及執行架構並不相同，本研究亦未在相同模型與執行預算下重現上述方法。因此，前述差異應視為根據各方法設計與本研究案例觀察提出的可能解釋，尚不能直接判定為效能差距的單一原因。針對 RQ1，本實驗顯示 CBR-APIKGAgent 在相同模型條件下改善了整體任務完成表現，但尚未達到 AppWorld 公開結果中的最佳水準。至於兩項核心機制各自的貢獻，將於下一節透過消融實驗進一步分析。

\section{實驗二：消融實驗與案例分析}

本節對應 RQ2、RQ3 與 RQ4。首先，透過消融實驗分別分析原子化經驗重用機制與動態局部修復機制對系統表現造成的影響，以回答 RQ2 與 RQ3。其次，本節延續消融實驗結果，透過可取得完整軌跡的成功與限制案例，觀察兩項機制如何在規劃與執行階段共同影響任務完成，以回答 RQ4。

\subsection{實驗組別}

消融實驗共包含四種設定，用以分別觀察原子化經驗與局部修復對任務完成結果的影響。四種設定皆於 AppWorld \texttt{test\_normal} 的 168 個任務進行評估，並採相同計分方式。除上述機制差異外，其餘模型、時間限制、重試上限與評估方式皆保持一致。

\subsection{消融實驗結果}

正式消融實驗以基礎系統作為比較基準，分別加入原子化經驗、動態局部修復，以及同時啟用兩項機制，以觀察各項機制對任務完成率的影響。四種設定皆使用相同任務集合，並以 TGC 與 SGC 作為比較指標，其中表中的 $\Delta$ 代表相較於 Baseline System 的提升幅度，單位為 pp。實驗結果顯示，僅啟用原子化經驗或僅啟用局部修復時，兩項指標皆高於基礎系統；同時啟用兩項機制的完整系統則取得四種設定中最高的 TGC 與 SGC。各組實際結果整理如表 \ref{tab:ablation_results} 所示。

\begin{table}[htbp]
    \centering
    \caption{消融實驗結果}
    \label{tab:ablation_results}
    \small
    \renewcommand{\arraystretch}{1.15}
    \setlength{\tabcolsep}{4pt}
    \begin{tabularx}{1\textwidth}{
        >{\raggedright\arraybackslash}X
        >{\centering\arraybackslash}p{0.13\textwidth}
        >{\centering\arraybackslash}p{0.13\textwidth}
        >{\centering\arraybackslash}p{0.13\textwidth}
        >{\centering\arraybackslash}p{0.13\textwidth}
    }
        \hline
        \textbf{Configuration} & \textbf{TGC (\%)} & \textbf{\ensuremath{\Delta}TGC (pp)} & \textbf{SGC (\%)} & \textbf{\ensuremath{\Delta}SGC (pp)} \\
        \hline
        Baseline System & 37.5 & --- & 14.3 & --- \\
        Atomic Experience Only & 43.5 & +6.0 & 25.0 & +10.7 \\
        Local Repair Only & 49.4 & +11.9 & 28.6 & +14.3 \\
        Full System & \textbf{51.8} & \textbf{+14.3} & \textbf{35.7} & \textbf{+21.4} \\
        \hline
    \end{tabularx}
\end{table}

為補充整體 TGC 與 SGC 指標，本研究進一步比較不同消融設定下的 task-level 評估結果差異。表 \ref{tab:task_level_ablation_comparison} 列出三類與研究問題相關的正向差集，用以觀察各機制可能帶來改善的任務集合。

\begin{table}[htbp]
    \centering
    \caption{消融設定之任務層級正向差集比較}
    \label{tab:task_level_ablation_comparison}
    \small
    \renewcommand{\arraystretch}{1.15}
    \begin{tabularx}{0.9\textwidth}{
        >{\raggedright\arraybackslash}X
        >{\centering\arraybackslash}p{0.22\textwidth}
    }
        \hline
        \textbf{Task-level outcome comparison} & \textbf{Number of tasks} \\
        \hline
        Baseline failed, AE-only succeeded & 19 \\
        Baseline failed, Local Repair-only succeeded & 32 \\
        Both single-mechanism settings failed, Full System succeeded & 8 \\
        \hline
    \end{tabularx}
\end{table}

各比較條件分別統計，前兩類可能包含相同任務。表 \ref{tab:task_level_ablation_comparison} 僅呈現不同消融設定下的任務評估結果差異，用以觀察各機制可能帶來改善的任務集合；兩項機制的實際作用方式仍需配合執行軌跡分析。因此，本表不應被解讀為完整的任務轉移矩陣，也不直接將任務成功歸因於單一機制。

後續將先分析原子化經驗與局部修復各自的影響，再透過完整系統與可取得完整軌跡的案例，討論兩項機制如何在規劃與執行階段共同影響任務完成。

\subsection{單一機制效益分析}

本節分別比較僅原子化經驗與僅局部修復兩種設定相較於基礎系統的差異，以回答 RQ2 與 RQ3。

\subsection*{原子化經驗之影響}

由表 \ref{tab:ablation_results} 可知，基礎系統之 TGC 為 37.5\%，SGC 為 14.3\%。僅原子化經驗設定之 TGC 為 43.5\%，SGC 為 25.0\%，相較於基礎系統分別提升 6.0 與 10.7 個百分點。此結果表示，在未啟用動態局部修復的情況下，原子化經驗設定仍呈現較高的任務層級與情境層級完成率。

進一步從任務層級比較來看，表 \ref{tab:task_level_ablation_comparison} 顯示「基礎系統失敗、僅原子化經驗成功」的任務共有 19 題。此差集表示，在部分任務中，加入原子化經驗後可觀察到原本未通過評估的任務轉為通過；然而，這些任務的改善情形仍需配合執行軌跡檢視，不能僅由差集結果直接歸因於單一機制。

以下以一個 Venmo payment request 修正任務作為代表性案例，說明其中一種可觀察到的改善情形。此案例用於呈現原子化經驗可能如何協助系統在規劃階段建立目標集合與篩選流程，並不代表所有僅原子化經驗成功的任務皆具有相同原因。

\noindent\textbf{範例：}
\begin{itemize}
    \item \textbf{案例任務：}使用者需修正前一晚向 roommates 發出的 Venmo payment requests，刪除原本金額少 10 美元的請求，並以相同描述與收款對象建立修正後的新請求。
    \item \textbf{消融差異：}此任務在基礎系統中未通過評估，但在僅原子化經驗設定中成功。
    \item \textbf{基礎系統觀察：}基礎系統將 roommate 身分視為缺失資訊，反覆要求使用者提供 roommates 的 names 或 email addresses，因而未能從既有資料中建立可用的目標集合。
    \item \textbf{機制作用：}僅原子化經驗設定的軌跡中，系統在初始規劃與後續重新規劃過程中，檢索到與寫入前目標集合建立、既有資料比對、目標篩選，以及使用 stable identifiers 維持操作對象一致性相關的策略經驗。此類經驗可能與後續規劃中的目標集合建立行為相關；在該軌跡中，系統將 roommates 這類隱含對象條件轉換為可查詢的 contact relationship 條件，先透過 phone contacts 查詢並以 roommate 關係取得 roommates 的 email；接著登入 Venmo、取得 sent payment requests，依 roommates email 與前一晚時間條件篩選待修正 requests，最後刪除原請求並建立金額增加 10 美元的新請求。
    \item \textbf{影響：}此代表性軌跡顯示，原子化經驗可能有助於在規劃階段建立目標集合與篩選流程，特別是將任務中的隱含對象條件轉換為可執行的資料取得步驟。不過，此案例不代表所有僅原子化經驗改善的任務皆具有相同原因。
\end{itemize}
\subsection*{動態局部修復之影響}

僅局部修復設定之 TGC 為 49.4\%，SGC 為 28.6\%，相較於基礎系統分別提升 11.9 與 14.3 個百分點。此結果表示，在不加入原子化經驗的情況下，執行失敗後的局部計畫修正仍可能使部分原本失敗的任務恢復並完成。進一步從任務層級比較來看，表 \ref{tab:task_level_ablation_comparison} 顯示「基礎系統失敗、僅局部修復成功」的任務共有 32 題。此差集可用於觀察局部修復可能帶來改善的任務集合，但各任務的實際改善方式仍需搭配執行軌跡分析，不能直接概括為同一種失敗原因。

從已檢視的代表性軌跡來看，局部修復可觀察到對不同類型執行問題的補救行為，例如授權缺失、必要前置資料不足、API 選擇不適合或查詢條件不精確等。這些例子顯示，動態局部修復可能與執行過程中利用局部錯誤證據修正當前失敗點的行為相關，使系統有機會在不重新規劃整個任務的情況下恢復執行。

\noindent\textbf{範例：}
\begin{itemize}
    \item \textbf{案例任務：}使用者需更新 Spotify 播放清單，依據兄弟姊妹在手機訊息中的建議調整歌曲內容。
    \item \textbf{消融差異：}此任務在基礎系統中失敗，但在僅局部修復設定中成功。
    \item \textbf{基礎系統觀察：}基礎系統能確認 sibling 關係，查到三位 sibling contacts，並取得包含歌曲新增與移除建議的 text messages。然而，後續流程步驟在辨識目標 Spotify playlist 時，未能將欲辨識的 roadtrip playlist 對應到實際的 Road Trip 播放清單；由於該步驟在未找到相符播放清單後即進入停止流程，後續歌曲搜尋、新增與移除操作皆未執行，因此 evaluation tests 中沒有產生預期的 playlist song 變更。
    \item \textbf{局部修復軌跡：}僅局部修復設定的初始流程曾使用 family 關係查詢 contacts 以尋找 family members 或 siblings，但結果為空，導致後續從空結果中擷取 sibling phone numbers 的步驟無法繼續。系統因此觸發局部修復，新增 phone contacts 查詢作為修復步驟，並改用 sibling 關係查詢。
    \item \textbf{機制作用：}在該軌跡中，修復步驟取得三位 sibling contacts 後，系統繼續搜尋 text messages、解析歌曲新增與移除建議、取得 playlist、搜尋歌曲，最後透過 Spotify playlist API 執行歌曲新增與移除，並通過 evaluation tests。
    \item \textbf{影響：}此代表性案例顯示，動態局部修復可能有助於在執行流程中根據錯誤回饋補足缺失的前置資料，協助任務從局部失敗點恢復執行。不過，此案例中的局部修復點與基礎系統的失敗點並不相同，因此只能說明一種可觀察到的修復行為，不能推論所有僅局部修復成功的任務皆具有相同原因。
\end{itemize}

整體而言，原子化經驗與動態局部修復在單獨啟用時皆高於基礎系統。針對 RQ2，僅原子化經驗設定使 TGC 與 SGC 分別提升 6.0 與 10.7 個百分點，代表性案例則顯示其可能透過提供資料取得、條件篩選與操作順序等策略參考，改善部分任務的規劃完整性。針對 RQ3，僅局部修復設定使 TGC 與 SGC 分別提升 11.9 與 14.3 個百分點，代表性案例則顯示其可能利用執行階段的局部錯誤證據修正失敗步驟，提升部分任務的恢復能力。在本組實驗中，僅局部修復設定之兩項指標皆高於僅原子化經驗設定；代表性案例顯示，兩項機制可分別在規劃與執行恢復階段產生作用，並對任務完成結果產生不同層次的影響。

\subsection{完整系統與代表性案例分析}
\label{sec:ablation_case_analysis}

本節對應 RQ4，目的在於說明原子化經驗與動態局部修復可能如何共同影響任務完成。相較於前一節以消融數據分析兩項機制各自對任務完成率的影響，本節先透過可取得完整軌跡的成功案例，觀察兩者在規劃與執行階段可能形成的互補關係；再透過限制案例，討論目前完整系統仍可能面臨的限制。

完整系統同時啟用原子化經驗與動態局部修復，其 TGC 為 51.8\%，SGC 為 35.7\%，為四種消融設定中最高。相較於基礎系統，完整系統之 TGC 提升 14.3 個百分點，SGC 提升 21.4 個百分點；相較於僅原子化經驗設定，完整系統之 TGC 與 SGC 分別提升 8.3 與 10.7 個百分點；相較於僅局部修復設定，完整系統之 TGC 與 SGC 分別提升 2.4 與 7.1 個百分點。

從任務層級差集來看，表 \ref{tab:task_level_ablation_comparison} 顯示「僅原子化經驗設定與僅局部修復設定皆失敗，但完整系統成功」的任務共有 8 題。此結果表示，部分任務僅在兩項機制同時啟用時通過評估，可作為觀察兩項機制可能互補的起點；但此差集本身不能直接證明兩項機制在所有任務中皆具有協同作用，仍需搭配具體執行軌跡分析。

\subsection*{互補案例分析}

根據可取得完整軌跡的完整系統成功案例可觀察到，原子化經驗可能在規劃階段提供目標集合建立與篩選策略的參考；若執行過程仍出現授權或前置資訊缺失，動態局部修復則可能依據當下錯誤補足必要步驟。此類案例顯示，兩項機制可能與不同階段的行為變化相關，並在部分任務中共同影響任務完成。

\noindent\textbf{範例：}
\begin{itemize}
    \item \textbf{案例任務：}使用者需將手機聯絡人中的 coworkers 與 friends 加入 Venmo friends，但僅限於目前尚未是 Venmo friends 的對象。
    \item \textbf{消融差異：}此任務在僅原子化經驗設定與僅局部修復設定中皆未通過評估，但在完整系統中成功。
    \item \textbf{單一機制觀察：}僅原子化經驗設定雖執行至新增好友步驟，但在將手機聯絡人對應到 Venmo 使用者時，搜尋結果包含了不屬於原始 coworkers 或 friends 的額外對象；後續篩選又未充分扣回「手機聯絡人中的 coworkers 或 friends，且尚未是 Venmo friends」這一任務條件，因此最後新增 13 位使用者，新增集合超出評估預期。僅局部修復設定則可處理 Venmo 授權錯誤與部分漏處理的查詢結果，但修復後仍沿用已被擴大的候選集合，因此最後新增 5 位使用者，仍未符合評估條件。
    \item \textbf{完整系統軌跡：}完整系統先從 phone contacts 取得候選對象，篩選 relationships 為 coworker 或 friend 的 contacts，形成候選好友 email 集合。接著系統嘗試透過 Venmo friends 查詢取得既有 Venmo friends，但遇到 401 authentication error；動態局部修復觸發後，新增 Venmo login 作為修復步驟，取得 Venmo access token 後回到原本的 friends 查詢流程。
    \item \textbf{機制作用：}在該軌跡中，完整系統於修復後抽取既有 Venmo friend emails，再與候選好友 email 集合做差集，形成待新增 email 集合，最後只新增應加入的 7 位使用者。
    \item \textbf{影響：}此代表性案例顯示，原子化經驗相關流程可能有助於規劃階段的目標集合建立與既有狀態比對，動態局部修復則可能在執行階段處理 Venmo 授權錯誤，協助原流程繼續執行。不過，此案例僅說明部分完整系統成功任務中可觀察到的互補關係，不能概括為所有任務皆具有相同原因。
\end{itemize}
\subsection*{完整系統的限制分析}

雖然完整系統在整體 TGC 與 SGC 上皆為四種設定中最高，但在個別任務中，並不保證一定優於僅啟用單一機制的設定。這類情況不宜直接解讀為原子化經驗與動態局部修復彼此干擾；較合理的解釋是，完整系統的效果仍會受到任務結構、搜尋式 API 語意，以及後續目標集合維持方式影響。

以下案例與前述成功案例同樣涉及 Venmo friends 與 phone contacts 的對應，但呈現的是另一種情況：完整系統雖能處理部分執行錯誤，仍可能在目標集合維持上遺漏應新增對象。此例僅呈現其中一種失敗型態，並非所有完整系統失敗案例的共同原因。

\noindent\textbf{範例：}
\begin{itemize}
    \item \textbf{案例任務：}使用者需將 Venmo 上的朋友名單調整為與手機聯絡人中的朋友一致，必要時執行新增與移除操作。
    \item \textbf{消融差異：}此任務中，基礎系統與僅局部修復設定皆成功，但完整系統失敗。
    \item \textbf{完整系統觀察：}完整系統雖在初始計畫中納入手機聯絡人中 friend 關係的篩選，但後續使用 venmo.search\_users 將 phone friends 對應到 Venmo users 時，未能完整保留所有應新增對象。結果中實際新增的 Venmo friends 少於評估預期，表示應新增對象未被完整納入操作目標集合。
    \item \textbf{局部修復限制：}局部修復較可能處理由執行結果或語意判斷所辨識的步驟失敗，例如登入、授權或缺少查詢步驟等問題；但當問題發生在上游目標集合建立或候選對象維持時，後續修復即使被觸發，也未必能完整恢復原始任務條件。
    \item \textbf{影響：}此例顯示，完整系統雖能結合經驗檢索與局部修復，但在涉及搜尋式 API 與寫入操作的任務中，仍需要明確維持「原始任務條件 → 候選結果 → 實際操作目標」之間的對應關係。若應新增對象在候選集合建立過程中被遺漏，後續局部修復即使能補足授權或查詢步驟，也未必能使操作目標集合恢復完整。此限制案例不代表完整系統失敗皆源於相同原因。
\end{itemize}
從系統流程來看，API 呼叫完成後並非僅依據是否發生執行錯誤判斷步驟結果，系統亦會利用任務描述、步驟目標、API 參數、回傳結果與可取得的相依資料進行語意判斷。然而，局部修復與重新規劃仍以目前步驟先被判定為失敗為前提。若搜尋 API 回傳的候選集合未完整保留原始任務條件，或原始任務限制未完整呈現於後續步驟及相依資訊中，該步驟仍可能被判定為成功。此時，若後續流程亦未產生足以暴露偏差的執行回饋，不完整或不一致的操作目標可能持續傳遞，而不一定能透過後續局部修復或重新規劃完全修正。

因此，針對 RQ4，消融結果與代表性案例可觀察到，原子化經驗多作用於規劃階段的任務結構與資料流策略，動態局部修復則多作用於執行失敗後的局部計畫調整；兩者因作用階段不同，可能在部分任務中形成互補，並與完整系統取得四種設定中的最佳整體結果相互呼應。然而，兩項機制的結合並非在所有任務中皆能帶來正向效果；當檢索經驗未能維持任務限制，或上游目標集合已產生偏差時，局部修復未必能恢復正確的任務條件。此結果提示未來可從三個方向持續改善：在經驗使用時納入任務限制、細化局部修復與重新規劃的適用範圍，以及建立更可靠的經驗累積機制，以維持任務條件、候選結果與實際操作目標之間的對應關係。

\section{實驗效度限制}

為降低實驗結果受到非研究變因影響，本研究於消融實驗中使用相同的任務集合、基礎模型、評估指標與主要執行設定，並僅改變原子化經驗與動態局部修復兩項機制是否啟用，以比較各機制對任務完成率的影響。然而，大型語言模型仍具有生成不確定性，即使在相同設定下，單次執行結果仍可能受到模型輸出差異、API 呼叫順序或中間推理結果影響。因此，本研究結果可反映各設定在相同控制條件下的相對表現；未來若進一步進行多次重複執行，則可更細緻估計模型輸出變動對結果的影響。

此外，本研究主要於 AppWorld \texttt{test\_normal} 上進行評估。AppWorld 提供標準化的多應用程式 API 環境與程式化 evaluation tests，使不同方法能在相同任務定義下比較任務完成情形。然而，真實 API 環境可能包含更動態的資料狀態、權限限制、文件不一致、非同步操作與外部服務錯誤等因素。因此，本研究結果能說明 CBR-APIKGAgent 在標準化 API 任務基準中的表現與限制，但其效果是否能完全推廣至真實 API 服務或其他任務環境，仍需進一步驗證。

\end{ZhChapter}

% 章節六 Chapter 6
\begin{ZhChapter}

\chapter{結論與未來研究}

\section{結論}

在此總結本研究的主要貢獻與實驗結果。

\section{未來研究}

在此說明可延伸的後續研究方向。

\end{ZhChapter}

% 新增你自己的章節... Add chapters here

% 參考文獻，請在段落上隨意註解，你同時需要 reference.bib
\clearpage
\phantomsection
\chapter*{參考文獻}
\addcontentsline{toc}{chapter}{參考文獻}

\printbibliography[heading=none]

% 附錄 Appendix
\chapter*{附錄 A 系統提示詞}
\phantomsection
\addcontentsline{toc}{chapter}{附錄 A 系統提示詞}

\setcounter{section}{0}
\renewcommand{\thesection}{A.\arabic{section}}
\renewcommand{\thesubsection}{\thesection.\arabic{subsection}}
\renewcommand{\theHsection}{appendix.A.\arabic{section}}
\renewcommand{\theHsubsection}{appendix.A.\arabic{section}.\arabic{subsection}}

本附錄整理 CBR-APIKGAgent 於任務規劃、局部修復與重新規劃階段所使用之主要提示詞。正文第四章僅說明各提示詞的輸入資訊與設計目的；完整提示內容則於本附錄列出，以利讀者對照系統實作。

\section{初始規劃提示詞}
\label{app:initial-planning-prompt}

【待補】整理初始規劃階段使用之提示詞，包含使用者任務、候選 API、規劃導向經驗與計畫格式限制。

\section{錯誤分類與可行回饋提示詞}
\label{app:error-feedback-prompt}

【待補】整理步驟失敗後用於產生暫定錯誤分類與可行回饋之提示詞，包含失敗步驟、失敗 API 資訊、原始錯誤訊息、相依資料脈絡與結構化失敗資訊。

錯誤分類提示詞中使用之分類集合依模組概念整理如下：

\begin{itemize}
    \item \textbf{記憶相關：} \texttt{oversimplification}、\texttt{false\_memory}、\texttt{retrieval\_failure}。
    \item \textbf{反思相關：} \texttt{progress\_misassessment}、\texttt{outcome\_misinterpretation}、\texttt{wrong\_causal\_attribution}、\texttt{reflection\_hallucination}。
    \item \textbf{規劃相關：} \texttt{constraint\_ignorance}、\texttt{impossible\_plan}、\texttt{inefficient\_plan}。
    \item \textbf{行動相關：} \texttt{plan\_misalignment}、\texttt{format\_error}、\texttt{parameter\_error}。
    \item \textbf{系統相關：} \texttt{step\_limit\_reached}、\texttt{tool\_execution\_error}、\texttt{llm\_limitation}、\texttt{environment\_error}。
    \item \textbf{驗證與其他：} \texttt{auth\_error}、\texttt{unknown}。
\end{itemize}

\section{局部修復策略選擇提示詞}
\label{app:repair-strategy-prompt}

【待補】整理局部修復階段用於選擇修復 API、相依資料使用方式與插入、替換、刪除策略之提示詞。

\section{重新規劃提示詞}
\label{app:replanning-prompt}

【待補】整理重新規劃階段使用之提示詞，包含原始任務、既有計畫、失敗資訊、已完成步驟結果，以及規劃導向與修復導向經驗。

\chapter*{附錄 B 程式碼}
\phantomsection
\addcontentsline{toc}{chapter}{附錄 B 程式碼}

\setcounter{section}{0}
\renewcommand{\thesection}{B.\arabic{section}}
\renewcommand{\thesubsection}{\thesection.\arabic{subsection}}
\renewcommand{\theHsection}{appendix.B.\arabic{section}}
\renewcommand{\theHsubsection}{appendix.B.\arabic{section}.\arabic{subsection}}

在此整理重要程式碼片段與關鍵實作說明。


\end{document}
